\documentclass{article}

% ============================================================
% PACKAGES
% ============================================================
\usepackage[utf8]{inputenc}
\usepackage[T1]{fontenc}

\usepackage[a4paper,right=1.4cm,top=1.4cm,left=1.4cm,bottom=1.8cm]{geometry}

\usepackage{amsmath,amssymb,amsthm,mathtools,bm}
\usepackage{booktabs,array}
\usepackage[ruled,vlined,linesnumbered]{algorithm2e}
\usepackage[inline]{enumitem}
\usepackage{xcolor}
\usepackage[colorlinks=true,linkcolor=blue!60!black,
 citecolor=green!50!black,urlcolor=blue!70!black]{hyperref}
\usepackage[expansion=false]{microtype}
\emergencystretch=1.5em

\usepackage[numbers,sort&compress]{natbib}

% ------------------------------------------------------------
% Theorem environments
% ------------------------------------------------------------
\theoremstyle{definition}
\newtheorem{definition}{Definition}[section]
\theoremstyle{plain}
\newtheorem{theorem}{Theorem}[section]
\newtheorem{proposition}[theorem]{Proposition}
\newtheorem{conjecture}[theorem]{Conjecture}
\newtheorem{hypothesis}[theorem]{Hypothesis}
\theoremstyle{remark}
\newtheorem{remark}{Remark}[section]

% ------------------------------------------------------------
% Convenience macros
% ------------------------------------------------------------
\newcommand{\R}{\mathbb{R}}
\newcommand{\eps}{\varepsilon}
\newcommand{\Sdist}{S}
\newcommand{\seclabel}[1]{\label{sec:#1}}
\newcommand{\secref}[1]{Section~\ref{sec:#1}}
\DeclareMathOperator*{\argmin}{argmin}

\title{Does Privacy Distort Alignment?}
\author{}
\date{}

% ============================================================
% DOCUMENT
% ============================================================
\begin{document}
\maketitle

\section{Thesis}\seclabel{thesis}
\paragraph{Setup.}
Let $\theta_0$ be a small instruction-tuned open-weight LM. A privacy
mechanism $m \in \mathcal{M}$ applied at budget $\eps$ acts either on the
adaptation data (text sanitization, token perturbation) or on the adaptation
procedure itself (DP-SGD), producing a perturbed model or
inference pipeline $\theta_{m,\eps}$. Each run yields two measurements:
a \emph{semantic distortion} score $\Sdist \in [0,1]$, computed on the
perturbed text or on generations relative to the unperturbed reference and, by
construction, \emph{without reference to} $\eps$; and an \emph{alignment
degradation} score $Y$, the drop relative to the no-privacy baseline on a fixed
evaluation suite (truthfulness, hallucination; \secref{metrics}).

\paragraph{Claim.}
The project's central research question: is DP-induced alignment degradation
driven by semantic distortion of the training/inference signal, or by
non-semantic side effects of DP noise? We formalize this as a simple linear model:
\begin{equation}\label{eq:model}
 Y_{i} = a_{m_i} + b_{m_i}\, d_i + c\, \Sdist_i + e_i,
\end{equation}
where $d_i$ is the \emph{privacy-level rank} of run $i$ within its own mechanism (the
four $\eps$ levels of M2/M3 and the four $p$ levels of M1, coded $1$--$4$ and
z-scored), $a_m$ and $b_m$ are a per-mechanism intercept and level slope, a
random intercept per base model is added to $a_m$, and $e_i$ is the error
term. $Y$, $d$, and $\Sdist$ are z-scored before fitting, so $b_m$ and $c$
are standardized effects on a common scale and can be compared directly.
Because $Y$ is defined as the drop relative to the same-model M0 baseline,
the M0 runs have $Y \equiv 0$ by construction: they define the zero point and
are not regression rows, so the regression uses the $216$ mechanism runs
($108$ per corpus).

\paragraph{Computing $d_i$.}
Each mechanism has exactly four settings, ordered from strictest to loosest:
M2 uses word-level $\eps \in \{1, 3, 6, 12\}$, M3 uses example-level
$\eps \in \{1, 4, 8, 16\}$, and M1 uses perturbation rates
$p \in \{0.50, 0.30, 0.15, 0.05\}$ (a higher rate is a stricter setting, so
the order is reversed). Within its own mechanism, a run at the strictest
setting gets $d_i = 1$, the next $d_i = 2$, up to $d_i = 4$ at the loosest;
the four values are then z-scored (giving $-1.16, -0.39, 0.39, 1.16$). Only
the \emph{order} of the settings is used, never their numeric values: this
is deliberate, because the numeric $\eps$ values are not comparable across
mechanisms (Remark~\ref{rem:eps}) and, within a mechanism, we claim only
that stricter means noisier---not any particular functional form. Rank
coding assumes exactly that and nothing more.

\paragraph{One regression, not three.}
Although $a_{m}$ and $b_{m}$ carry a mechanism index, model~\eqref{eq:model}
is a single ordinary linear regression, estimated jointly on all $216$ rows
via indicator (dummy) variables. Written out for mechanisms $m \in \{1,2,3\}$:
\begin{equation}\label{eq:dummies}
  Y_i \;=\; \sum_{m=1}^{3} a_m\,\pmb{1}[m_i{=}m]
        \;+\; \sum_{m=1}^{3} b_m\, d_i\,\pmb{1}[m_i{=}m]
        \;+\; c\, \Sdist_i \;+\; e_i,
\end{equation}
a linear model with seven coefficients and a $216 \times 7$ design matrix,
fitted by least squares in one pass (in \texttt{statsmodels}:
\texttt{Y \textasciitilde{} C(mechanism) + C(mechanism):d + S}, with the
per-model random intercept added via \texttt{MixedLM}). Each row activates
the intercept and slope of its own mechanism only, which is why the model
behaves as if fitted per mechanism---but the $\Sdist$ column carries no
indicator, so $c$ is shared by all rows, and every mechanism's runs
contribute evidence about the same $c$. This is the design's point: the
mechanisms disagree about how much distortion a given privacy level buys,
but they pool into one estimate of what distortion costs. Joint estimation
also makes Hypothesis~\ref{hyp:dom} directly testable: $\hat c$ and each
$\hat b_m$ come from one fit with their joint covariance, so
$|c| > |b_m|$ is a Wald test on linear combinations of coefficients, not a
comparison across separate regressions.

\begin{hypothesis}[Pre-registered, primary: minimum effect]\label{hyp:main}
$|c| \ge \Delta$ in model~\eqref{eq:model}, with $\Delta = 0.10$
standardized units fixed in advance. In words: semantic distortion has a
practically meaningful effect on alignment degradation, not merely a nonzero
one. In the mediation decomposition of the total effect of $\eps$ on $Y$, the
indirect path through $\Sdist$ carries a non-trivial share.
\end{hypothesis}

\begin{hypothesis}[Secondary: dominance]\label{hyp:dom}
$|c| > |b_{m}|$ for each mechanism $m$: semantic distortion
predicts degradation more strongly than the nominal privacy level does.
(The level slope $b_m$ is mechanism-specific by construction, so the
comparison is a Wald test per mechanism; it may be underpowered and is
therefore secondary.)
\end{hypothesis}

\begin{hypothesis}[Secondary: moderation by entity density]\label{hyp:mod}
Distortion hurts alignment more when the text is fact-dense. Formally: in
model~\eqref{eq:model} extended with a $\Sdist \times$ dataset interaction
term, the coefficient of $\Sdist$ is larger on the entity-rich corpus
(Nemotron-PII) than on the entity-poor one (Alpaca). If confirmed, this
bounds how far results from generic instruction data generalize to truly
private corpora---and implies mechanisms should be evaluated on entity-rich
text to be trusted.
\end{hypothesis}

\begin{conjecture}[Directional, secondary]\label{conj:shape}
The privacy--alignment relationship is not monotone in $\eps$ alone: at matched
nominal privacy levels, mechanisms with low $\Sdist$ (semantics-preserving
sanitization) degrade alignment
strictly less than mechanisms with high $\Sdist$ (random token replacement),
and DP-SGD occupies an intermediate
regime in which degradation is dominated by the \emph{direct} (optimization /
calibration) path, i.e.\ large $\hat b$-attributable residual at small
$\Sdist$.
\end{conjecture}

\paragraph{Identification.}
Within a single mechanism, $\Sdist$ and $\eps$ are strongly correlated by
construction: more noise means more distortion. So a regression on one
mechanism alone cannot separate $b_m$ from $c$. We need variation in
$\Sdist$ at fixed $\eps$. It comes from two sources:
\begin{enumerate}[left=0pt,itemsep=1pt,topsep=1pt]
 \item \textbf{Cross-mechanism contrast.} Different mechanism families reach
 a similar privacy level with very different amounts of semantic damage.
 This is the main source of identification, and the reason the full
 mechanism grid (\secref{design}) is required.
 \item \textbf{A non-DP distortion control.} Random token
 replacement/deletion, tuned to the same distortion levels, produces
 $\Sdist$ without any privacy mechanism. If this arm degrades alignment
 as much as the DP arms at equal $\Sdist$, then semantic damage---not DP
 itself---is the cause. If not, the difference points to DP-specific,
 non-semantic effects.
\end{enumerate}
We check collinearity with variance inflation factors (VIF) in every fitted
model. If VIF $>10$ for $\Sdist$, we switch to a simpler comparison:
mechanisms against each other at each fixed privacy level. That comparison
does not suffer from $\eps$--$\Sdist$ confounding at all.

\begin{remark}[On matching $\eps$ across mechanism types]\label{rem:eps}
Nominal budgets \emph{cannot be compared} across mechanism families: DP-SGD
gives example-level $(\eps,\delta)$-DP, while SanText-style sanitization gives
word-level metric LDP. We do not pretend they are on the same scale. In
model~\eqref{eq:model}, the privacy level enters only through the privacy-level rank
$d$, coded within its own mechanism and given a per-mechanism slope $b_m$.
\emph{$\Sdist$ is the common axis across mechanisms.} If alignment degradation is predicted better
by the measurable, comparable quantity ($\Sdist$) than by the incomparable
nominal one ($\eps$), that is a useful result on its own.
\end{remark}

\begin{remark}[The null is a result]\label{rem:null}
The hypothesis and $\Delta$ are pre-registered: committed to the repository
before the main experimental pass. Three outcomes are possible. If the
confidence interval for $c$ lies entirely outside $(-\Delta, \Delta)$,
Hypothesis~\ref{hyp:main} is confirmed. If it lies entirely inside, a two
one-sided equivalence test (TOST) lets us state that the semantic pathway is
practically absent---not just ``not detected''. In that case degradation
comes from non-semantic pathways (optimization instability, calibration
shifts), so semantics-preserving mechanism design would \emph{not} help
alignment. This is a valid negative result, and it directly informs the
mitigation deliverable. If the interval straddles $\Delta$, the result is
inconclusive and is reported as such. All three outcomes lead to a written
conclusion.
\end{remark}

\section{Experimental design}\seclabel{design}
\paragraph{Mechanism grid ($4$ arms).}
\begin{enumerate}[left=0pt,itemsep=1pt,topsep=1pt]
 \item \textbf{M0 --- No privacy.} Standard LoRA
 fine-tuning on clean data; defines the zero point of $Y$ and $\Sdist$.
 \item \textbf{M1 --- Random token perturbation (non-DP control).} Uniform
 token replacement/deletion at rates
 $p \in \{0.05, 0.15, 0.30, 0.50\}$ chosen to span the $\Sdist$ range
 realized by the DP arms.
 \item \textbf{M2 --- Semantic text sanitization (metric LDP).}
 SanText-style word-level sanitization with
 embedding-space exponential mechanism (TEM as a
 variant if time permits), $\eps \in \{1, 3, 6, 12\}$ (word-level).
 \item \textbf{M3 --- DP fine-tuning.} DP-SGD applied to LoRA adapters:
 per-example gradient clipping plus Gaussian noise, implemented directly
 on the adapter parameters via \texttt{torch.func} (\texttt{vmap} over
 per-example gradients is cheap because the adapter parameter count is
 small); $\eps \in \{1, 4, 8, 16\}$ at $\delta = 10^{-5}$
 (example-level).
\end{enumerate}

\paragraph{Models ($3$ families, dense only).}
Qwen3-1.7B, Llama-3.2-1B(-Instruct), and Falcon3-1B-Instruct, with the 3--4B siblings as
a scale check in week~3 if compute allows. We deliberately exclude the newest
Qwen3.5 small series from the critical path: its hybrid Gated-DeltaNet~+
sparse-MoE architecture makes per-example gradient computation awkward, and
debugging privacy accounting on a novel architecture is too risky for
week~2. Qwen3.5 stays as an optional extra arm for the
perturbation-only mechanisms (M1--M2), which do not depend on architecture.

\paragraph{Grid and compute.}
$3$ models $\times$ $4$ mechanisms $\times$ $4$ perturbation levels $\times$ $3$ seeds
$= 144$ runs ($108$ excluding M0, which has no privacy level axis; M0 runs
$3\times3=9$ seeds/models, total $117$ training runs plus evaluation passes).
Each run is a LoRA fine-tune of a $\le$2B model on $\sim$10k instruction
pairs: $\approx$20--40 GPU-minutes on a single A100/H100, i.e.\ the full grid
is $\sim$60--80 GPU-hours---feasible within the school's allocation, with the
pilot ($1$ model $\times$ all arms $\times$ $1$ seed $=13$ runs) costing under
one GPU-day.

\paragraph{Data.}
The privacy mechanisms need a text corpus to act on: M1--M2 rewrite the
training text, and M3 fine-tunes on it privately. This corpus plays the role
of the ``sensitive data''. The two corpora should differ in \emph{entity
density}, because privacy risk concentrates in names, numbers, and facts, and
mechanisms treat these tokens differently from common words---rare tokens sit
in sparse regions of embedding space, so sanitization replaces them with
distant, fact-changing substitutes, and DP-SGD suppresses exactly the
memorization of such rare sequences. The entity-density contrast lets us test
whether the distortion--alignment link depends on it (\secref{analysis}).

\emph{Entity-rich corpus: Nemotron-PII}
(\texttt{nvidia/Nemotron-PII}, CC~BY~4.0): $100$k synthetic records
(emails, clinical notes, forms, free text) across $50$+ industries, with
\emph{gold span-level annotations} for $55$+ PII/PHI entity types. It is
built precisely as a stand-in for sensitive data: realistic PII with zero
re-identification risk. Two properties matter for us. First, the gold spans
make $\Sdist_{\mathrm{ent}}$ (\secref{metrics}) \emph{exact}---entity
preservation is checked against annotations instead of a noisy NER model.
Second, because every entity is known, we can directly measure whether
fine-tuned models reproduce or corrupt specific entities, which links the
distortion measurements to memorization. One adaptation step is required:
the records are documents, not instruction pairs, so we wrap each record
into a templated task (``Summarize this note'', ``Extract the appointment
date'', ``Draft a reply to this email''), a standard instructionization
step performed once and reused by all arms.

\emph{Low-entity-density corpus: Alpaca} (a fixed 10k subset of its 52k
instruction--response pairs): generic assistant tasks with few named
entities and numbers, permissively licensed, and the standard small
instruction-tuning corpus, which keeps results comparable to prior work.
Evaluation data are disjoint from adaptation data by construction.

\section{Metrics}\seclabel{metrics}
\paragraph{Semantic distortion $\Sdist$ (fixed before the main pass).}
\emph{Primary:} mean cosine distance between sentence
embeddings of original vs.\ perturbed/sanitized
training text (for M1--M2) or of generations from $\theta_{m,\eps}$ vs.\
$\theta_0$ on a fixed probe set (for M3, where no perturbed text exists).
\emph{Robustness:} (i) semantic entropy of
generations, connecting $\Sdist$ to the hallucination literature at the level
of meaning rather than surface tokens; (ii) NLI-based fact-level agreement
(AlignScore). The primary measure is fixed now so that it cannot be changed
after seeing the results; the robustness measures appear only in a
sensitivity table.

\emph{Decomposition of $\Sdist$.} $\Sdist$ is additionally reported as two
components: $\Sdist_{\mathrm{ent}}$, the fraction of named entities and
numbers not preserved between original and perturbed text (checked against
gold spans on Nemotron-PII; NER-based match on the second corpus), and
$\Sdist_{\mathrm{gen}}$, the residual embedding drift after masking
entities. Prediction: $\Sdist_{\mathrm{ent}}$ drives hallucination and
truthfulness degradation, while $\Sdist_{\mathrm{gen}}$ drives
instruction-following degradation. Fitting model~\eqref{eq:model} with the
two components in place of $\Sdist$ tests this directly.

\paragraph{Alignment degradation $Y$.}
Three components, all judge-free (no LLM judge anywhere in the
measurement chain): \textbf{TruthfulQA} (MC1/MC2, truthfulness),
\textbf{IFEval} (instruction following; each prompt carries programmatically
verifiable constraints, so scoring is exact), and an \textbf{entity-fidelity
probe} (hallucination): the model is asked to reproduce facts from held-out
Nemotron-PII records, and answers are checked against the gold entity
spans---correct, corrupted, or fabricated---with no judge noise. TruthfulQA
and IFEval are subsampled to $300$--$500$ items so a full evaluation pass per
checkpoint takes minutes. $Y$ is the standardized drop relative to the
same-model M0 baseline, averaged across the three components (per-component
coefficients reported in the appendix).
\emph{Secondary probes of the non-semantic (direct) pathway:} expected
calibration error on TruthfulQA MC, and training-loss variance across seeds
as a proxy for optimization instability.

\section{Analysis plan}\seclabel{analysis}
Fit model~\eqref{eq:model} by restricted maximum likelihood
(\texttt{statsmodels} MixedLM), $144$ observations, random intercepts per
model family; cluster-robust standard errors by (model $\times$ mechanism)
cell. Report:
\begin{enumerate}[left=0pt,itemsep=1pt,topsep=1pt]
 \item $\hat c$ with $95\%$ CI; the pre-registered minimum-effect test
 of Hypothesis~\ref{hyp:main} against $\Delta = 0.10$, and the TOST
 equivalence test for the null branch (Remark~\ref{rem:null}); per-mechanism
 Wald tests of $|c| > |b_{m}|$ (Hypothesis~\ref{hyp:dom}).
 \item The mediation decomposition: proportion of the total privacy-level effect
 (within each mechanism) carried by the $\Sdist$ path, with
 bootstrap CIs over seeds.
 \item Privacy--distortion--alignment curves: $Y$ vs.\ privacy level and $Y$ vs.\
 $\Sdist$, per mechanism and model, across both datasets---the
 curator's deliverable~3.
 \item The entity-density interaction (Hypothesis~\ref{hyp:mod}): refit
 model~\eqref{eq:model} with a $\Sdist \times$ dataset term; a
 significant interaction means the distortion--alignment link depends
 on how fact-dense the corpus is.
 \item The $\Sdist$ decomposition fit: $\Sdist_{\mathrm{ent}}$ and
 $\Sdist_{\mathrm{gen}}$ in place of $\Sdist$ (\secref{metrics}),
 with per-benchmark coefficients.
 \item Diagnostics: VIF for $\Sdist$, residuals-by-mechanism, and the
 within-$\eps$-stratum contrast as the collinearity fallback.
\end{enumerate}

\paragraph{Reproducibility.}
Fixed seeds alone do not determinize GPU training. The repository therefore
pins a full environment lockfile,
enables \texttt{torch.use\_deterministic\_algorithms(True)} where kernels
permit, logs every config, seed, checkpoint hash, and metric to
\texttt{wandb}. Most importantly, the $3$ seeds per cell give the
within-condition variance: every plotted point has an error bar, and no claim
rests on a single run. The pre-registration (hypothesis,
$\Delta$, primary $\Sdist$ metric, exclusion rules) is a timestamped commit
made before the week-3 pass.

\section{Benchmarks}\seclabel{benchmarks}
\begin{table}[h]
\centering\small
\begin{tabular}{@{}llll@{}}
\toprule
Component & Role & Primary metric & Guards against \\
\midrule
TruthfulQA (MC1/MC2) & truthfulness $Y$ & accuracy drop vs.\ M0 & imitative falsehoods \\
IFEval & instruction following $Y$ & constraint pass-rate drop & judge-based scoring \\
Entity-fidelity probe & hallucination $Y$ & gold-span match rate & judge noise \\
Embedding distance & primary $\Sdist$ & cosine distance & metric cherry-picking \\
Semantic entropy / AlignScore & robustness $\Sdist$ & rank agreement w/ primary & measure-specific artifacts \\
M1 random perturbation & non-DP control & $Y$ at matched $\Sdist$ & ``DP is special'' confound \\
ECE + loss variance & direct-path probes & calibration / stability & unexplained residual \\
\bottomrule
\end{tabular}
\caption{Evaluation suite and the role of each component.}
\label{tab:bench}
\end{table}

\section{Milestones}\seclabel{milestones}
\begin{table}[h]
\centering\small
\begin{tabular}{@{}p{0.06\columnwidth}p{0.86\columnwidth}@{}}
\toprule
Wk & Milestone \\
\midrule
M1 & Setup: repo skeleton, environment lockfile, eval harness over TruthfulQA
 + IFEval subsamples and the entity-fidelity probe, $\Sdist$ metrics implemented and unit-tested on
 synthetic perturbations; literature map. \textbf{Gate:} clean M0 baseline
 numbers on all three models. \\
M2 & Mechanisms: M1--M3 implemented; \textbf{DP-SGD/LoRA per-example-gradient
 smoke test on all three architectures} (the largest technical risk, so it
 is done early on purpose); pilot $=1$ model $\times$ all arms $\times 1$
 seed. \textbf{Gate:} the pilot must show clearly different $\Sdist$ values
 at matched privacy level across mechanisms; pre-registration commit. \\
M3 & Main pass: full $144$-run grid across both datasets;
 privacy--distortion--alignment curves; fit model~\eqref{eq:model};
 mediation decomposition; diagnostics. \\
M4 & Mitigation probe (one semantics-preserving variant, e.g.\ TEM or
 sanitization with entity protection, evaluated at the worst-case cell);
 equivalence analysis if null; paper + code release. \\
\bottomrule
\end{tabular}
\caption{Milestones. The week-2 gate is the go/no-go point: if DP-SGD fails on
an architecture, that model is replaced from the dense fallback pool
(Falcon3-3B, Llama-3.2-3B) rather than debugged past the gate.}
\label{tab:miles}
\end{table}

\end{document}
